\chapter{Enough terminal to be dangerous}
\label{ch:terminal}

\chapterbrief
  {The dozen commands this book actually uses.}
  {Paths, permissions, and where you are.}
  {Confidence that you cannot break anything by accident.}

Skip this chapter if \code{cd}, \code{ls}, \code{chmod +x} and \code{sudo} are
already reflexes. Otherwise it is fifteen minutes that will save you hours.

\section{Where am I?}

The shell always has a \emph{current directory}. Almost every confusing error in
setup work comes from being somewhere you did not think you were.

\begin{shell}
pwd                 # print working directory -- where am I?
ls                  # list files here
ls -la              # ... including hidden ones, with details
cd ~/projects       # change directory
cd ..               # up one level
cd -                # back to where I just was
\end{shell}

\begin{concept}{The three special paths}
\begin{description}[leftmargin=0pt, style=nextline]
  \item[\code{\textasciitilde}] your home directory. On the Pi that is
        \code{/home/martin}; on macOS \code{/Users/martin}.
  \item[\code{.}] here, the current directory. \code{./script.sh} means ``the
        script in \emph{this} folder'' --- as opposed to a command found on your
        \code{PATH}.
  \item[\code{..}] one level up.
\end{description}

A path starting with \code{/} is \emph{absolute} --- it means the same thing from
anywhere. Anything else is \emph{relative} to where you currently are, which is
why \code{pwd} is the first thing to check when a command cannot find a file.
\end{concept}

\section{Looking at files}

\begin{shell}
cat file.txt              # dump the whole thing
less file.txt             # page through it (q to quit, / to search)
head -20 file.txt         # first 20 lines
tail -20 file.txt         # last 20 lines
tail -f logfile           # follow it live -- Ctrl-C to stop
grep "error" file.txt     # lines containing "error"
grep -rn "TODO" src/      # search a whole tree, with line numbers
\end{shell}

\code{tail -f} is how you watch a program's log while it runs. You will use it in
Volume One when the kiosk misbehaves.

\section{Editing on the Pi}

You need one editor that works over SSH. \code{nano} is the least painful:

\begin{shell}
nano ~/.bashrc
\end{shell}

Arrow keys move. \textbf{Ctrl-O then Enter} saves. \textbf{Ctrl-X} exits. Those
two shortcuts are all you need, and they are printed at the bottom of the screen
--- where \code{\^{}O} means Ctrl-O.

\begin{trap}{If you find yourself in \texttt{vi}}
Some commands (notably \code{git commit} with no \code{-m}) open a full-screen
editor that ignores everything you type. That is \code{vim}, and it is modal.

To escape: press \textbf{Escape}, then type \textbf{\code{:q!}} and press Enter.
That quits without saving.

To make git use nano instead, permanently:
\begin{shell}
git config --global core.editor nano
\end{shell}
\end{trap}

\section{Permissions, and \texttt{chmod +x}}

Every file carries permissions: who may read it, write it, and \emph{execute} it.

\begin{shell}
ls -l scripts/deploy.sh
\end{shell}

\begin{trace}
-rwxr-xr-x  1 martin  staff  1204 Jul 26 14:02 scripts/deploy.sh
 ^^^
 owner may read, write, execute
\end{trace}

A shell script without the execute bit gives you \code{Permission denied} even
though you can plainly read it. Fix:

\begin{shell}
chmod +x scripts/deploy.sh
\end{shell}

\begin{trap}{Scripts lose the execute bit in transit}
Copy a script through a zip file, a Windows filesystem, or some cloud drives and
the execute bit is silently dropped.

If a script you \emph{know} is fine says \code{Permission denied}, run
\code{chmod +x} on it before suspecting anything else. Git does preserve the
bit, which is one more reason to use git.
\end{trap}

\section{\texttt{sudo}, and when not to use it}

\code{sudo} runs one command as the administrator. Installing packages and
editing system configuration need it. \textbf{Your own project files never do.}

\begin{trap}{Never \texttt{sudo} a build}
If you run \code{sudo mvn compile} even once, the files in \code{target/} end up
owned by root. Your normal user then cannot overwrite them, so the \emph{next}
build fails with a permission error --- and the obvious-looking fix is more
\code{sudo}, which digs the hole deeper.

If it has already happened:
\begin{shell}
sudo chown -R $USER:$USER ~/midnight-brewer
\end{shell}
\end{trap}

\section{Keys that save you}

\begin{center}\small
\begin{tabular}{@{}L{0.22\linewidth}L{0.70\linewidth}@{}}
\toprule
\thead{Key} & \thead{Does} \\
\midrule
\textbf{Tab} & Completes file and command names. Press it constantly. It is also how you check a path exists before pressing Enter. \\
\textbf{Up arrow} & Previous command. Edit and re-run rather than retyping. \\
\textbf{Ctrl-C} & Stop the running program. \\
\textbf{Ctrl-D} & End of input --- also logs you out of SSH. \\
\textbf{Ctrl-R} & Search your command history. Type a few letters of a command you ran last week. \\
\textbf{Ctrl-A / Ctrl-E} & Jump to start / end of the line. \\
\bottomrule
\end{tabular}
\end{center}

\begin{concept}{Tab completion is a correctness tool}
It is not just about typing speed. If Tab does not complete a filename, that
file \emph{does not exist} where you think it does --- and you have found a bug
before running anything.

Beginners type full paths and get ``No such file or directory''. Experienced
people press Tab, see nothing happen, and immediately run \code{ls}.
\end{concept}

\section{Environment variables}

A variable your shell passes to programs it runs. This book uses one:

\begin{shell}
export PI=martin@brewer.local     # set it for this session
echo $PI                          # read it back
\end{shell}

\code{export} only lasts until you close the terminal. To make it permanent, add
the line to your shell's startup file:

\begin{shell}
# which shell am I using?
echo $SHELL

# zsh (macOS default)
echo 'export PI=martin@brewer.local' >> ~/.zshrc

# bash (most Linux, WSL2)
echo 'export PI=martin@brewer.local' >> ~/.bashrc
\end{shell}

Then open a new terminal, or run \code{source \textasciitilde/.zshrc} to apply
it now.

\begin{checkpoint}{Prove the basics}
\begin{shell}
cd ~
pwd                       # prints your home directory
mkdir -p projects/test
cd projects/test
echo "hello" > file.txt
cat file.txt              # prints: hello
ls -l                     # shows file.txt with its permissions
cd ~
rm -rf ~/projects/test    # clean up
\end{shell}

If all of that worked and you understand what each line did, you know enough
terminal for this entire book.
\end{checkpoint}
